\section{Modelo do painel fotovoltaico}\label{sec:painel}

Arquivos: \arq{painel/celula_params.m}, \arq{kcl_modulo.m}, \arq{parametros_painel.m},
\arq{datasheet_painel.m}, \arq{curva_pv.m}. O painel \'e modelado como um circuito el\'etrico
equivalente.

\subsection{O circuito equivalente de 2 diodos}

Uma c\'elula solar se comporta como uma fonte de corrente (a luz gera el\'etrons) em paralelo com
dois diodos (perdas por recombina\c{c}\~ao), uma resist\^encia em paralelo $R_p$ (fuga) e uma
resist\^encia em s\'erie $R_s$ (contatos e metaliza\c{c}\~ao).

\begin{figure}[htbp]
\centering
\begin{circuitikz}[scale=0.9,transform shape]
  \draw (0,0) to[isource,l=$\Ipv$] (0,3) -- (2,3);
  \draw (2,3) to[D,l=$D_1$] (2,0);
  \draw (2,3) -- (4,3) to[D,l=$D_2$] (4,0);
  \draw (4,3) -- (6,3) to[R,l=$R_p$] (6,0);
  \draw (6,3) to[R,l=$R_s$] (8.6,3) to[short,-o] (9.6,3);
  \draw (0,0) -- (9.6,0) to[short,-o] (9.6,0);
  \draw[<->,>=Stealth] (10,0.1) -- node[right]{$V$} (10,2.9);
  \node[gray,font=\scriptsize] at (4.8,-0.6) {$V_d$ \'e a tens\~ao nos tr\^es ramos em paralelo; $V=V_d-R_sI$};
\end{circuitikz}
\caption{Circuito equivalente de uma c\'elula (ou de um grupo de c\'elulas em s\'erie) no modelo
de 2 diodos.}
\label{fig:circuito2d}
\end{figure}

A lei das correntes (KCL) d\'a a equa\c{c}\~ao do m\'odulo:
\begin{equation}\label{eq:2d}
I=\Ipv-I_{01}\!\left(e^{V_d/(a_1V_t)}-1\right)-I_{02}\!\left(e^{V_d/(a_2V_t)}-1\right)-\frac{V_d}{R_p},
\qquad V_d=V+R_sI .
\end{equation}

\begin{itemize}
\item $\Ipv$: fotocorrente, proporcional \`a luz; \'e o que o painel ``gera''.
\item $I_{01},I_{02}$: correntes de satura\c{c}\~ao. O diodo 1 representa a difus\~ao de
 portadores ($a_1=1$); o diodo 2, a recombina\c{c}\~ao na regi\~ao de dep\'osito ($a_2=1{,}2$).
\item $V_t=kT/q$: tens\~ao t\'ermica ($\approx25{,}7$\,mV por c\'elula a 25\,$^\circ$C). No m\'odulo
 multiplica-se pelo n\'umero de c\'elulas em s\'erie.
\item $V_d/R_p$: corrente de fuga pela resist\^encia paralela.
\end{itemize}

O c\'odigo devolve o \emph{res\'iduo} $F=0$ quando o ponto $(V,I)$ est\'a sobre a curva
(C\'odigo~\ref{lst:kcl}). Essa forma \emph{expl\'icita} evita resolver a equa\c{c}\~ao impl\'icita
$I(V)$ em todo lugar: basta testar se um ponto a satisfaz. Isso torna a extra\c{c}\~ao de
par\^ametros (Cap\'itulo~\ref{sec:extracao}) simples.

\codigointeiro{painel/kcl_modulo.m}{res\'iduo da equa\c{c}\~ao do m\'odulo e inclina\c{c}\~ao $dI/dV$.}{lst:kcl}

\begin{logica}
\begin{enumerate}
\item O m\'odulo todo \'e tratado como \emph{um} grupo de $N_s$ c\'elulas em s\'erie
 (\texttt{q.Ncel = p.Ns}). As grandezas dependentes de $G$ e $T$ v\^em de \arq{celula_params}.
\item \texttt{Vd = V + Rs\_s*I}: tens\~ao na jun\c{c}\~ao = tens\~ao do terminal mais a queda em $R_s$.
\item \texttt{F} \'e a eq.~\eqref{eq:2d} passada toda para um lado: fotocorrente menos os dois
 diodos menos a fuga menos a corrente de sa\'ida. Se $F=0$, o ponto est\'a na curva.
\item \texttt{g} \'e a condut\^ancia total da jun\c{c}\~ao, $g=-\partial I_c/\partial V_d$.
\item \texttt{dIdV}: por deriva\c{c}\~ao impl\'icita de $F(V,I)=0$,
 $\dfrac{dI}{dV}=-\dfrac{\partial F/\partial V}{\partial F/\partial I}=-\dfrac{g}{1+R_sg}$.
 Essa inclina\c{c}\~ao \'e usada na condi\c{c}\~ao $dP/dV=0$ da extra\c{c}\~ao.
\end{enumerate}
\end{logica}

\subsection{Depend\^encia com irradi\^ancia e temperatura}

\arq{celula_params(G,T,p)} concentra tudo o que muda com $G$ (W/m$^2$) e $T$ (K); a f\'isica fica
num \'unico lugar.

\codigointeiro{painel/celula_params.m}{grandezas do grupo de c\'elulas em fun\c{c}\~ao de $G$ e $T$.}{lst:celula}

\begin{center}\small
\begin{tabularx}{\textwidth}{@{}l>{\raggedright\arraybackslash}p{6.4cm}>{\raggedright\arraybackslash}X@{}}
\toprule
Grandeza & Lei usada & Intui\c{c}\~ao \\ \midrule
Fotocorrente & $\Ipv=\dfrac{G}{G_n}\big(I_{pv,n}+K_i(T-T_n)\big)$ & Mais luz, mais corrente; calor aumenta um pouco a corrente \\[6pt]
Satura\c{c}\~ao & $I_{0k}=I_{0k,n}\left(\dfrac{T}{T_n}\right)^{3}\exp\!\left[\dfrac{E_g}{k}\Big(\dfrac1{T_n}-\dfrac1T\Big)\right]$ & Diodos ``vazam'' mais a quente; por isso $V_{oc}$ cai com $T$ \\[8pt]
Tens\~ao t\'ermica & $V_t=kT/q$ & Cresce com $T$ \\
Resist\^encia paralela & $R_p(G)=R_{p,ref}\,G_n/G$ & Ver abaixo \\
Resist\^encia s\'erie & $N_{cel}R_s$ & C\'elulas em s\'erie somam resist\^encias \\
\bottomrule
\end{tabularx}
\end{center}

\begin{atencao}[: gap efetivo]
Na teoria da jun\c{c}\~ao, $E_g$ \'e o gap do sil\'icio (1,12\,eV). Com esse valor o modelo previa
$K_v=-0{,}106$\,V/K para o KC200GT, contra $-0{,}123$\,V/K do datasheet. Por isso $E_g$ virou um
par\^ametro \emph{efetivo}, ajustado pelo $K_v$ medido. Ele absorve efeitos que a lei simples
ignora.
\end{atencao}

\subsubsection{Por que $R_p$ depende da irradi\^ancia (D10)}

Com $R_p$ constante, a fuga $V/R_p$ quase n\~ao muda com a luz, mas a fotocorrente cai
proporcionalmente a $G$. Em 200\,W/m$^2$ a fuga vira uma fra\c{c}\~ao enorme da corrente gerada:
no ajuste do ZTP-160M ($R_p\approx41\,\Omega$) isso previa \textbf{26\,\% de perda de efici\^encia
em 200\,W/m$^2$}, contra 3 a 5\,\% t\'ipicos de m\'odulos monocristalinos. A lei de De Soto, Klein
e Beckman (2006) faz a fuga escalar com a luz. Em $G=G_n$ nada muda, ent\~ao a extra\c{c}\~ao
(toda em STC) n\~ao \'e afetada. Um subm\'odulo no escuro ($G=0$) d\'a $R_p=\infty$ e fuga zero,
sem NaN, pois $V_d/\infty=0$.

\subsection{Os dados de cat\'alogo}

\codigo{painel/datasheet_painel.m}{23}{41}{dados do ZTP-160M, com as convers\~oes de \%/K para A/K e V/K.}{lst:ds}

\begin{center}\small
\begin{tabular}{@{}ll@{}}
\toprule
Dado (ZTP-160M, datasheet v4.3) & Valor \\ \midrule
$I_{sc}$ / $V_{oc}$ & 9,02\,A / 22,64\,V \\
$I_{mp}$ / $V_{mp}$ & 8,20\,A / 19,57\,V \\
$P_{max}$ & 160\,W (mas $V_{mp}I_{mp}=160{,}47$\,W: o cat\'alogo arredonda) \\
$K_i$ & $+0{,}05$\,\%/K de $I_{sc}$ \\
$K_v$ & $-0{,}34$\,\%/K de $V_{oc}$ \\
$\gamma_P$ (coef.\ de $P_{max}$) & $-0{,}4$\,\%/K (usado s\'o para validar) \\
Arranjo & 33 c\'elulas, 3 subm\'odulos de 11 (hip\'otese: 1 bypass por coluna) \\
\'Area, massa, NMOT & 0,82\,m$^2$, 8,9\,kg, 45\,$^\circ$C \\
\bottomrule
\end{tabular}
\end{center}

O fator de idealidade do modelo de 1 diodo \'e $a=1{,}1$. Para o KC200GT, Villalva usa 1,3, mas o
ZTP-160M \'e monocristalino com fator de forma 0,78 e 1,3 n\~ao o alcan\c{c}a sem $R_s\le0$. O
valor 1,1 \'e o maior, em passos de 0,1, que ajusta.

\subsection{A curva P-V}

\arq{curva_pv(G,T,p,nPts)} varre a tens\~ao de 0 at\'e uma estimativa folgada de $V_{oc}$ e, a
cada tens\~ao, resolve a corrente do m\'odulo (com sombra, por \arq{painel_sombreado},
Cap\'itulo~\ref{sec:sombra}). Cada ponto usa a solu\c{c}\~ao do anterior como chute inicial
(\emph{warm start}), o que acelera o Newton. O resultado traz as curvas I-V e P-V, o m\'aximo
global $(V_{mp},P_{mp})$ e a tabela de todos os m\'aximos locais (\texttt{islocalmax}).

\codigo{painel/curva_pv.m}{12}{33}{varredura de tens\~ao com \emph{warm start} e detec\c{c}\~ao dos picos.}{lst:curvapv}
